\section{Direct Supersession Certificates}
\label{sec:method}

\subsection{Inputs and temporal eligibility}

Let $\mathcal C$ contain source-attributed claims.
A dated claim $c$ has text $x_c$, identifier $\mathrm{id}_c$, and finite timestamp $t_c$.
The timestamp identifies the supplied source state or recorded change time.
A publication date alone does not establish an unobserved effective date.
The query supplies a cutoff $\tau$ and a relevance ranking $\pi_q$ over the corpus.
Our output contains at most $k$ claims that remain eligible at $\tau$.

The pairwise gate $G(d,c)$ means that $d$ supplies a replacement for $c$.
It combines a shared-attribute decision with a changed-value decision.
The gate uses fixed claim features or a fixed pairwise procedure.
It need not be transitive, symmetric, or correct against world knowledge.
We study the evidence policy induced by this gate.
Interpreting its output as factual validity additionally requires correct replacement decisions and correctly ordered change times.
Concurrent values require a gate that distinguishes coexistence from replacement.

Define eligible evidence by
\begin{equation}
\begin{split}
A_G(\tau)=\{c:\;&t_c\leq\tau,\ \nexists d\in\mathcal C\text{ such that}\\
&t_c<t_d\leq\tau\text{ and }G(d,c)\}.
\end{split}
\label{eq:eligibility}
\end{equation}
Thus a claim remains eligible until an accepted later replacement becomes available at the cutoff.
Equal timestamps do not establish replacement order.
Undated claims receive no historical certificate.
The current-answer interface retains them as unresolved evidence, matching the archived system's policy.

\subsection{A certificate for every cutoff}

For each dated claim, define its first replacement boundary:
\begin{equation}
b_G(c)=\min\{t_d:t_d>t_c,\ G(d,c)\},
\label{eq:boundary}
\end{equation}
with $\min\varnothing=\infty$.
Store a witness $w_G(c)$ that attains this boundary.
The smallest identifier selects among witnesses with the same timestamp.
The certificate is
\begin{equation}
I_G(c)=[t_c,b_G(c)).
\label{eq:certificate}
\end{equation}
The interval determines eligibility.
The witness identifies the source responsible for its endpoint.
An infinite endpoint means that the corpus contains no accepted later replacement.

The selected context is
\begin{equation}
R_G(q,\tau)=\operatorname{Top}_k\langle c\in\pi_q:\tau\in I_G(c)\rangle.
\label{eq:context}
\end{equation}
Filtering precedes the context budget.
The operation preserves the relative relevance order of retained claims.
It does not fill empty slots with excluded claims.
Historical and current retrieval use the same certificates.
For dated claims, the final current mask retains exactly those with infinite endpoints.

Repeated values remain separate source records.
A return to an earlier value creates a new record with its own start time.
The earlier occurrence retains its original endpoint.
The method therefore represents recurrence without reopening an expired source record.
For the Northstar illustration, each appointment closes the preceding record through its own direct decision.

\subsection{Exact compilation}

\begin{figure}[t]
\centering
\fbox{\begin{minipage}{0.93\columnwidth}
\small
\textbf{Inputs:} dated claims, fixed gate $G$, complete postings.\\
\textbf{Output:} boundary $b(c)$ and witness $w(c)$.
\begin{enumerate}
\item Set $b(c)=\infty$ and $w(c)=\varnothing$ for every claim.
\item Initialize the unresolved index $U=\varnothing$.
\item Process timestamp batches $B$ in increasing order.
\item For each $d\in B$, process identifiers in order.
\item Find unresolved candidates $c$ from complete postings.
\item If $G(d,c)$, set $(b(c),w(c))=(t_d,d)$.
\item Remove each resolved $c$ from all postings in $U$.
\item After all comparisons for $B$, insert its claims into $U$.
\end{enumerate}
\end{minipage}}
\caption{Exact compilation. Steps 4--7 run inside each timestamp batch. Removing resolved claims avoids repeated work. Delayed batch insertion prevents replacement between simultaneous claims.}
\label{fig:compiler}
\end{figure}


Figure~\ref{fig:compiler} gives the compiler.
It processes timestamp batches in increasing order.
An inverted index stores unresolved older claims.
Each new claim queries complete candidate postings and evaluates the gate against unresolved candidates.
The first accepted replacement fixes an older claim's endpoint.
The compiler then removes that older claim from the postings.
No later replacement can improve its earliest endpoint.
Claims from the current batch enter the postings only after all comparisons finish.
This preserves equal-time claims and deterministic witness ties.

For the supplied text matcher, a positive frame threshold permits complete token blocking.
Accepted nonempty frames share a frame token.
All empty frames use one common posting key.
A nonpositive threshold uses one global posting key.
Custom gates use a full unresolved scan unless their candidate index has an equivalent completeness guarantee.
Blocking changes the comparisons examined, but not the certificates produced.

The supplied gate uses the earlier detector's fixed parsing rules.
It compares relation frames, named phrases, and changed numeric or textual values.
We preserve these rules throughout the compiler comparisons.
No gold subject, relation, answer value, or chronology enters the compiler.
Section~\ref{sec:protocol} identifies the annotations used only to construct questions and score results.

Let $n=|\mathcal C|$.
Let $L$ count posting entries, $V$ posting visits, and $P$ unique pair tests.
If one gate evaluation costs at most $g$, construction costs
\begin{equation}
O(n\log n+L+V+Pg)
\label{eq:cost}
\end{equation}
after feature extraction.
Output storage is $O(n)$ words.
Uninformative postings and long repeated-value runs can still require quadratic work.
We therefore report pair tests and posting visits separately from elapsed time.
